% Created 2016-02-23 mar. 09:35
\documentclass[a4paper,11pt,twoside,openright]{article}
\usepackage[utf8]{inputenc}
\usepackage[T1]{fontenc}
\usepackage{fixltx2e}
\usepackage{graphicx}
\usepackage{longtable}
\usepackage{float}
\usepackage{wrapfig}
\usepackage{rotating}
\usepackage[normalem]{ulem}
\usepackage{amsmath}
\usepackage{textcomp}
\usepackage{marvosym}
\usepackage{wasysym}
\usepackage{amssymb}
\usepackage{hyperref}
\tolerance=1000
\usepackage{a4wide}
\usepackage[T1]{fontenc}
\usepackage{mathptmx}
\usepackage[scaled=.90]{helvet}
\usepackage{courier}
\usepackage{fancyhdr}
\usepackage{tikz}
\author{Manuel Grizonnet (CNES), Julien Michel (CNES), Victor Poughon (CNES)}
\date{Décembre 2013}
\title{Formation à l'Orfeo ToolBox: Guide des Travaux Pratiques}
\hypersetup{
  pdfkeywords={},
  pdfsubject={},
  pdfcreator={Emacs 24.4.1 (Org mode 8.2.10)}}
\begin{document}

\maketitle
\tableofcontents

\pagestyle{fancy}
\fancyhf{}
\fancyhead[LE,RO]{\bfseries\thepage}
\fancyhead[LO]{\bfseries\rightmark}
\fancyhead[RE]{\bfseries\leftmark}
\fancyfoot[LE,RO]{\vspace{0pt}\includegraphics[height=20pt]{../../../../Artwork/logoVectoriel.png}}
\fancyfoot[LO,RE]{\vspace{0pt}\includegraphics[height=20pt]{Images/logo_cnes.png}}
\fancyfoot[C]{\vspace{2pt}\footnotesize{Formation Orfeo ToolBox - INRA 2016}}

\section{Avant-propos}
\label{sec-1}
\subsection{Données}
\label{sec-1-1}

\subsection{Les logiciels}
\label{sec-1-2}

Afin de réaliser les exercices, les logiciels suivants sont
nécessaires :

\begin{itemize}
\item \textbf{Orfeo ToolBox} 5.2 ou plus récent avec les applications
\item \textbf{Monteverdi} 3.0 ou plus récent
\item \textbf{QGis} 2.4 ou plus récent
\end{itemize}

Pour ce qui concerne l'installation de l'\textbf{Orfeo ToolBox} et
de  \textbf{Monteverdi}, vous pouvez consulter l' \href{http://www.orfeo-toolbox.org/packages/OTBSoftwareGuide.pdf}{Orfeo ToolBox Cookbook}.

Pour installer \textbf{QGis} consultez la documentation du logiciel qui est
disponible sur le  \href{http://www.qgis.org/}{site du projet}.

\newpage

\begin{center}
\vspace*{10cm}
\tiny Creative Commons Attribution-ShareAlike 3.0 Unported License\\
\normalsize
\includegraphics[width=0.1\textwidth]{Images/CC-licence.png}
\end{center}
\newpage

\section{Généralités}
\label{sec-2}
\subsection{Utiliser \textbf{Monteverdi} et \textbf{QGis}}
\label{sec-2-1}
\subsubsection{Description}
\label{sec-2-1-1}
\paragraph{Résumé}
\label{sec-2-1-1-1}

\paragraph{Pré-requis}
\label{sec-2-1-1-2}


\paragraph{Objectifs}
\label{sec-2-1-1-3}

\subsubsection{Étapes}
\label{sec-2-1-2}

\subsubsection{Solutions}
\label{sec-2-1-3}
\subsection{le mécanisme des applications \textbf{Orfeo ToolBox}}
\label{sec-2-2}
\subsubsection{Description}
\label{sec-2-2-1}
\paragraph{Résumé}
\label{sec-2-2-1-1}

\paragraph{Pré-requis}
\label{sec-2-2-1-2}


\paragraph{Objectifs}
\label{sec-2-2-1-3}

\subsubsection{Étapes}
\label{sec-2-2-2}

\subsubsection{Solutions}
\label{sec-2-2-3}
\subsection{Les mécanismes internes de l'\textbf{Orfeo ToolBox}}
\label{sec-2-3}
\subsubsection{Description}
\label{sec-2-3-1}
\paragraph{Résumé}
\label{sec-2-3-1-1}

\paragraph{Pré-requis}
\label{sec-2-3-1-2}


\paragraph{Objectifs}
\label{sec-2-3-1-3}

\subsubsection{Étapes}
\label{sec-2-3-2}

\subsubsection{Solutions}
\label{sec-2-3-3}
\section{Imagerie THR optique, des pré-traitements au SIG}
\label{sec-3}
\subsection{Pré-traitement de l'imagerie THR optique}
\label{sec-3-1}
\subsubsection{Description}
\label{sec-3-1-1}
\paragraph{Résumé}
\label{sec-3-1-1-1}

\paragraph{Pré-requis}
\label{sec-3-1-1-2}


\paragraph{Objectifs}
\label{sec-3-1-1-3}

\subsubsection{Étapes}
\label{sec-3-1-2}

\subsubsection{Solutions}
\label{sec-3-1-3}
\subsection{Segmentation et export vers un SIG}
\label{sec-3-2}
\subsubsection{Description}
\label{sec-3-2-1}
\paragraph{Résumé}
\label{sec-3-2-1-1}

\paragraph{Pré-requis}
\label{sec-3-2-1-2}


\paragraph{Objectifs}
\label{sec-3-2-1-3}

\subsubsection{Étapes}
\label{sec-3-2-2}

\subsubsection{Solutions}
\label{sec-3-2-3}
\section{Classification supervisée pour les séries multi-temporelles (Landsat 8)}
\label{sec-4}
\subsection{Réaliser une classification supervisée}
\label{sec-4-1}
\subsubsection{Description}
\label{sec-4-1-1}
\paragraph{Résumé}
\label{sec-4-1-1-1}

\paragraph{Pré-requis}
\label{sec-4-1-1-2}


\paragraph{Objectifs}
\label{sec-4-1-1-3}

\subsubsection{Étapes}
\label{sec-4-1-2}

\subsubsection{Solutions}
\label{sec-4-1-3}
\subsection{Evaluer, nettoyer et valoriser une carte de classification}
\label{sec-4-2}
\subsubsection{Description}
\label{sec-4-2-1}
\paragraph{Résumé}
\label{sec-4-2-1-1}

\paragraph{Pré-requis}
\label{sec-4-2-1-2}


\paragraph{Objectifs}
\label{sec-4-2-1-3}

\subsubsection{Étapes}
\label{sec-4-2-2}

\subsubsection{Solutions}
\label{sec-4-2-3}
\section{Traitements SAR pour l'imagerie Sentinel 1}
\label{sec-5}
\subsection{Calibration, débruitage et extraction de primitives pour les images SAR}
\label{sec-5-1}
\subsubsection{Description}
\label{sec-5-1-1}
\paragraph{Résumé}
\label{sec-5-1-1-1}

\paragraph{Pré-requis}
\label{sec-5-1-1-2}


\paragraph{Objectifs}
\label{sec-5-1-1-3}

\subsubsection{Étapes}
\label{sec-5-1-2}

\subsubsection{Solutions}
\label{sec-5-1-3}
% Emacs 24.4.1 (Org mode 8.2.10)
\end{document}
